% 来源：sources/2016-2025-期中-甲-历年真题汇编.pdf 第 10 页
\begin{problem}{15 分}{10cm}
设有下列 $n$ 阶行列式
\[
D_n=\begin{vmatrix}
7&-6&0&\cdots&0&0\\
-1&7&-6&0&\cdots&0\\
0&-1&7&-6&0&\cdots\\
\vdots&\vdots&\vdots&\vdots&\vdots&\vdots\\
0&\cdots&0&-1&7&-6\\
0&0&\cdots&0&-1&7
\end{vmatrix}.
\]
\begin{enumerate}[label=(\arabic*),leftmargin=2.2em,itemsep=0.25em,topsep=0.3em]
  \item 计算 $D_1,D_2$；
  \item 假设当 $n\le k$ 时有 $D_n=\dfrac{6^{n+1}-1}{5}$，证明：$D_{k+1}=\dfrac{6^{k+2}-1}{5}$。
\end{enumerate}
\end{problem}

\begin{problem}{10 分}{7.5cm}
设 $B=\begin{pmatrix}1&1&1\\0&1&1\\0&0&1\end{pmatrix}$，矩阵 $A$ 满足方程 $BA=B^{T}B$，求 $|A^{*}-2E|$。
\end{problem}

\begin{problem}{10 分}{7cm}
设 $n$ 阶实方阵 $A$ 满足 $A^{2}=A$，证明：$r(A)+r(A-E)=n$。
\end{problem}

\begin{problem}{20 分}{12cm}
设有四元线性方程组
\[
\begin{cases}
x_1+x_2-x_3+2x_4=0,\\
2x_1+3x_2-3x_3+5x_4=0,\\
3x_1+x_2-2x_3+3x_4=0,\\
2x_1+x_2-5x_3-x_4=0.
\end{cases}
\]
\begin{enumerate}[label=(\arabic*),leftmargin=2.2em,itemsep=0.25em,topsep=0.3em]
  \item 写出该方程组的系数矩阵 $A$；
  \item 利用初等行变换将矩阵 $A$ 化为阶梯形矩阵 $U$；
  \item 求 $r(A)$；
  \item 写出该方程组的通解。
\end{enumerate}
\end{problem}

\begin{problem}{10 分}{7.5cm}
设实方阵 $A$ 的伴随矩阵为 $A^{*}=\begin{pmatrix}1&0&0\\0&1&0\\1&0&1\end{pmatrix}$，且 $|A|>0$，已知矩阵 $B$ 满足 $AB=E+3A$，求矩阵 $B$。
\end{problem}

\begin{problem}{25 分}{12cm}
设矩阵 $A=\begin{pmatrix}3&1&0\\0&3&1\\0&0&3\end{pmatrix}$。求
\begin{enumerate}[label=(\arabic*),leftmargin=2.2em,itemsep=0.25em,topsep=0.3em]
  \item $A^{2}$；
  \item $A^{3}$；
  \item $A^{100}$；
  \item $A^{-1}$；
  \item 一个三次实系数多项式 $f(x)=x^{3}+ax^{2}+bx+c$ 使得 $f(A)$ 为三阶零方阵。
\end{enumerate}
\end{problem}

\begin{problem}{10 分}{7.5cm}
设有 $n$ 阶实方阵 $A=(a_{ij})$ 满足：当 $i=j$ 时，$a_{ij}=2021$；当 $i<j$ 时，$a_{ij}=i^{j}$；当 $i>j$ 时，$a_{ij}=(i+j)!$。证明：$|A|\ne 0$。
\end{problem}

\begin{problem}{10 分}{7.5cm}
设 $A$ 为 $n$ 阶方阵，且 $r(A)=r>0$。证明存在秩为 $r$ 的实方阵 $B$ 和 $C$，使得 $AB=CA$。
\end{problem}
